% !TeX root = ../../Book5.tex
% 纲要标题：### 33.4 Picard 函子、层化与代表对象
% 纲要位置：几何学计划纲要.md，第 6445 行。
% 正文状态：未撰写；本文件目前是框架。

% 纲要内容开始
% 
% * 相对线丛与来自基底的线丛之商
% * fppf 层化及局部线丛的下降障碍
% * 刚化、基点与 Poincaré 线丛
% * 一般可表性主版本及其代数空间结论
% 
%   明确外引如下版本：$f:X\to S$ 平坦、固有、有限呈示，且任意基变换 $T\to S$ 后 $\mathcal O_T\to(f_T)_*\mathcal O_{X_T}$ 为同构，则相对 Picard fppf 层由代数空间表示。是否为概形及是否存在未刚化通用线丛，需要分别给出附加条件。[R12]
% 
% * 代表空间、Picard 叠与通用线丛存在性的区别
% 
% 纲要内容结束

\section{Picard 函子、层化与代表对象}
\label{g5:sec:3304}

待撰写

\subsection{相对线丛与来自基底的线丛之商}
\label{g5:subsec:3304:01}
待撰写

\subsection{fppf 层化及局部线丛的下降障碍}
\label{g5:subsec:3304:02}
待撰写

\subsection{刚化、基点与 Poincaré 线丛}
\label{g5:subsec:3304:03}
待撰写

\subsection{一般可表性主版本及其代数空间结论}
\label{g5:subsec:3304:04}
待撰写

\subsection{代表空间、Picard 叠与通用线丛存在性的区别}
\label{g5:subsec:3304:05}
待撰写
% 纲要条目排为编号小节标题，并进入目录与书签；说明段仍保留于纲要注释中。
